\documentclass[10pt,a4paper]{amsart}
\usepackage[margin=19mm]{geometry}
\usepackage{amsmath,amssymb,mathtools}
\usepackage[scr=rsfs,cal=euler]{mathalfa}
\usepackage{xfrac}
\usepackage[unicode,hidelinks,bookmarks=false]{hyperref}
\newcommand{\ie}{i.e.\ }
\newcommand{\cf}{cf.\ }
\newcommand{\resp}{respectively\ }
\newcommand{\Subsection}{\subsection}
\providecommand{\nobreakdash}{\nobreak\hbox{-}}
\setlength{\emergencystretch}{2em}
\multlinegap0pt
\usepackage{amssymb}
\usepackage{xcolor}
\newtheorem{theorem}{Theorem}
\newtheorem{proposition}{Proposition}
\newtheorem{lemma}{Lemma}
\title{Weighted and Truncated Tail Index Estimation under Random Censoring: A Unified Full-Range Framework}
\author{Abdelhakim Necir, Nour Elhouda Guesmia, Djamel Meraghni}
\date{Source: arXiv:2605.13650v1 (CC BY 4.0)}
\begin{document}
\maketitle

\noindent\emph{Source and licence.} Weighted and Truncated Tail Index Estimation under Random Censoring: A Unified Full-Range Framework. Version arXiv:2605.13650v1 (CC BY 4.0), distributed by arXiv under the Creative Commons Attribution 4.0 International licence, http://creativecommons.org/licenses/by/4.0/, as recorded in the arXiv licence metadata for this exact version and shown on the abstract page https://arxiv.org/abs/2605.13650v1. Full licence notice: \texttt{licenses/arxiv-2605.13650-CC-BY-4.0.txt}.\par\smallskip

\noindent\textit{Editorial scope.} Counted unit: the article's proposition prop:linearization (source0.tex lines 2568--2588). The article's complete original proof of it is delivered with it (source0.tex lines 2590--2711, one \texttt{proof} environment of 122 lines). No mathematics is written, repaired or re-derived: the statement and the proof are the article's own lines, in the article's own notation, and no retained line is substituted or re-flowed.  Retained uncounted as the source-local context the proof needs: lines 144--147; lines 149--152; lines 470--473; lines 854--873; lines 2330--2381; lines 2493--2508.  Nothing was omitted from any retained span, so the package carries no omission record.  No retained line contains a citation, so the wrapper carries no bibliography and no bibliography entry.  No retained line contains a figure environment, an image-inclusion command or drawing code, so no image is delivered and the package declares no asset dependency.  Editorial formatting only, in addition: an AMS \texttt{amsart} preamble that restates the article's own declarations and macros used by the retained lines, namely the environments \texttt{theorem}, \texttt{proposition}, \texttt{lemma} on separate counters, together with the standard packages those lines need.  The retained environments print in this document as Theorem 1, Proposition 1, Lemma 1 in order of appearance, whereas the article prints the same items with the numbering of its own full document.  The complete native source of the article as distributed in the arXiv e-print for this exact version is bundled unchanged as \texttt{sources/200525/source0.tex}.
\begin{equation}
\lim_{t\rightarrow\infty} \frac{\overline{F}(tx)}{\overline{F}(t)} =
x^{-1/\gamma_{1}}, \label{RVF}%
\end{equation}

\begin{equation}
\lim_{t\rightarrow\infty} \frac{\overline{G}(tx)}{\overline{G}(t)} =
x^{-1/\gamma_{2}}, \label{RVG}%
\end{equation}

\begin{equation}
a_{n}\rightarrow0,\qquad na_{n}\rightarrow\infty,\qquad a_{n}=o(k/n).
\label{assump_an}%
\end{equation}

\begin{theorem}
[Consistency]\label{theorem2}

Assume that the survival functions $\overline{F}$ and $\overline{G}$ satisfy
the first-order regular variation conditions \eqref{RVF} and \eqref{RVG},
respectively. Let $k=k_{n}$ be an intermediate sequence such that
\[
k\to\infty\qquad\text{and} \qquad k/n\to0, \quad\text{as } n\to\infty.
\]
Let $(a_{n})$ be a positive sequence satisfying \eqref{assump_an}, and define
\[
m_{n} := \lfloor n a_{n} \rfloor.
\]
Then, for every $\beta>1$,
\[
\widehat{\gamma}_{1,k}^{(\mathrm{NA,tr})}(\beta) \overset{\mathbf{P}%
}{\longrightarrow} \gamma_{1}, \qquad\text{as } n\to\infty.
\]

\end{theorem}

\begin{proposition}
\label{prop:app-functional} Set
\[
Z_{k}:=Z_{n-k:n}, \qquad T_{n,a}:=\frac{u_{n}}{Z_{k}}, \qquad\alpha_{n}%
:=\frac{\beta}{\widehat{p}_{k}},
\]
and, for $1\leq x\leq T_{n,a}$, define
\[
Q_{n}(x) := \frac{\overline{F}_{n}^{(\mathrm{NA})}(xZ_{k})} {\overline{F}%
_{n}^{(\mathrm{NA})}(Z_{k})}, \qquad R_{n,a}(x) := \frac{\overline{F}%
_{n}^{(\mathrm{NA},a)}(xZ_{k})} {\overline{F}_{n}^{(\mathrm{NA})}(Z_{k})},
\]
and
\[
c_{n} := \frac{\overline{F}_{n}^{(\mathrm{NA})}(u_{n})} {\overline{F}%
_{n}^{(\mathrm{NA})}(Z_{k})}.
\]
Then, for every $\beta>1$,
\begin{equation}
\widehat{\gamma}_{1,k}^{(\mathrm{NA,tr})}(\beta) = \alpha_{n} \int%
_{1}^{T_{n,a}} x^{-1} \left(  Q_{n}(x)\right)  ^{\alpha_{n}} \,dx - \alpha
_{n}(\log T_{n,a})\,c_{n}^{\alpha_{n}}. \label{eq:app-functional-Q}%
\end{equation}
Moreover, for $1\leq x\leq T_{n,a}$,
\[
Q_{n}(x)=R_{n,a}(x)+c_{n},
\]
so that
\begin{equation}
\widehat{\gamma}_{1,k}^{(\mathrm{NA,tr})}(\beta) = \alpha_{n} \int%
_{1}^{T_{n,a}} x^{-1} \left(  R_{n,a}(x)+c_{n}\right)  ^{\alpha_{n}} \,dx -
\alpha_{n}(\log T_{n,a})\,c_{n}^{\alpha_{n}}. \label{eq:app-functional-R}%
\end{equation}
Finally, using
\[
D_{n,a}^{(\beta)}(x) = \sqrt{k}\,x^{-\beta/\gamma} \left(  R_{n,a}%
(x)-x^{-1/\gamma_{1}} \right)  , \qquad1\leq x\leq T_{n,a},
\]
that is,
\[
R_{n,a}(x) = x^{-1/\gamma_{1}} + \frac{1}{\sqrt{k}}x^{\beta/\gamma}%
D_{n,a}^{(\beta)}(x),
\]
we obtain the exact functional representation
\begin{equation}
\widehat{\gamma}_{1,k}^{(\mathrm{NA,tr})}(\beta) = \alpha_{n} \int%
_{1}^{T_{n,a}} x^{-1} \left(  x^{-1/\gamma_{1}} + c_{n} + \frac{1}{\sqrt{k}%
}x^{\beta/\gamma}D_{n,a}^{(\beta)}(x) \right)  ^{\alpha_{n}} dx - \alpha
_{n}(\log T_{n,a})\,c_{n}^{\alpha_{n}}. \label{eq:app-functional-D}%
\end{equation}

\end{proposition}

\begin{lemma}
\label{lem:boundary-truncation} Assume that $(a_{n})$ satisfies
\[
a_{n}\to0,\qquad na_{n}\to\infty,\qquad a_{n}=o(k/n),
\]
and, in addition,
\[
\sqrt{k} \left(  \frac{na_{n}}{k}\right)  ^{\beta} \log\!\left(  \frac
{k}{na_{n}}\right)  \to0.
\]
Then
\[
\sqrt{k}\,\alpha_{n}(\log T_{n,a})c_{n}^{\alpha_{n}} = o_{\mathbf{P}}(1).
\]

\end{lemma}

\begin{proposition}
[Linearization of the estimator]\label{prop:linearization} Assume that the
conditions of Theorem~\ref{theorem2} hold. Assume, in addition, that
\[
\sqrt{k} \left(  \frac{na_{n}}{k}\right)  ^{\beta} \log\!\left(  \frac
{k}{na_{n}}\right)  \to0.
\]
Then, for every $\beta>1$, we have
\[
\sqrt{k}\left(  \widehat{\gamma}_{1,k}^{(\mathrm{NA,tr})}(\beta) - \gamma_{1}
\right)  = \left(  \frac{\beta}{p}\right)  ^{2} \int_{1}^{T_{n,a}}
x^{-1+1/\gamma_{1}} D_{n,a}^{(\beta)}(x)\,dx + B_{n} + o_{\mathbf{P}}(1),
\]
where
\[
B_{n} := \sqrt{k} \left[  \frac{\beta}{p} \int_{1}^{T_{n,a}} x^{-1} \left(
\frac{\overline{F}(Z_{k}x)}{\overline{F}(Z_{k})} \right)  ^{\beta/p} \,dx -
\gamma_{1} \right]  .
\]

\end{proposition}

\begin{proof}
We start from the exact functional representation established in
Proposition~\ref{prop:app-functional}:
\[
\widehat{\gamma}_{1,k}^{(\mathrm{NA,tr})}(\beta)=\alpha_{n}\int_{1}^{T_{n,a}%
}x^{-1}Q_{n}(x)^{\alpha_{n}}\,dx-\alpha_{n}(\log T_{n,a})c_{n}^{\alpha_{n}},
\]
where
\[
\alpha_{n}:=\frac{\beta}{\widehat{p}_{k}}.
\]
Since $\widehat{p}_{k}\overset{\mathbf{P}}{\longrightarrow}p$, we have
\[
\alpha_{n}\overset{\mathbf{P}}{\longrightarrow}\frac{\beta}{p}.
\]
Set
\[
A_{n}(x):=\frac{\overline{F}(Z_{k}x)}{\overline{F}(Z_{k})}.
\]
We decompose
\[
\begin{aligned}
		\sqrt{k}\left(
		\widehat{\gamma}_{1,k}^{(\mathrm{NA,tr})}(\beta)-\gamma_{1}
		\right)
		&=
		\sqrt{k}
		\left[
		\alpha_n\int_{1}^{T_{n,a}}x^{-1}Q_n(x)^{\alpha_n}\,dx
		-
		\frac{\beta}{p}\int_{1}^{T_{n,a}}x^{-1}A_n(x)^{\beta/p}\,dx
		\right]  \\
		&\quad+
		\sqrt{k}
		\left[
		\frac{\beta}{p}\int_{1}^{T_{n,a}}x^{-1}A_n(x)^{\beta/p}\,dx
		-\gamma_1
		\right]  \\
		&\quad-
		\sqrt{k}\,\alpha_n(\log T_{n,a})c_n^{\alpha_n}.
	\end{aligned}
\]
The second term is exactly $B_{n}$. By Lemma~\ref{lem:boundary-truncation},
the last term is $o_{\mathbf{P}}(1)$.

It remains to linearize the first term. By the uniform stochastic expansion of
the Nelson--Aalen tail ratio,
\[
Q_{n}(x)=A_{n}(x)+\{Q_{n}(x)-A_{n}(x)\},
\]
and Taylor's formula applied to the map $u\mapsto u^{\beta/p}$ gives
\[
Q_{n}(x)^{\beta/p}=A_{n}(x)^{\beta/p}+\frac{\beta}{p}A_{n}(x)^{\beta
/p-1}\{Q_{n}(x)-A_{n}(x)\}+r_{n}(x),
\]
where
\[
\sqrt{k}\int_{1}^{T_{n,a}}x^{-1}|r_{n}(x)|\,dx=o_{\mathbf{P}}(1).
\]
Moreover, the replacement of $\alpha_{n}$ by $\beta/p$ in the preceding
integral contributes only $o_{\mathbf{P}}(1)$ after multiplication by
$\sqrt{k}$. Hence
\[
\begin{aligned}
		&\sqrt{k}
		\left[
		\alpha_n\int_{1}^{T_{n,a}}x^{-1}Q_n(x)^{\alpha_n}\,dx
		-
		\frac{\beta}{p}\int_{1}^{T_{n,a}}x^{-1}A_n(x)^{\beta/p}\,dx
		\right] \\
		&\qquad =
		\left(\frac{\beta}{p}\right)^{2}
		\int_{1}^{T_{n,a}}x^{-1}A_n(x)^{\beta/p-1}
		\sqrt{k}\{Q_n(x)-A_n(x)\}\,dx
		+
		o_{\mathbf P}(1).
	\end{aligned}
\]
From the definition of the weighted truncated process,
\[
D_{n,a}^{(\beta)}(x)=\sqrt{k}\,x^{-\beta/\gamma}\left(  R_{n,a}%
(x)-x^{-1/\gamma_{1}}\right)  ,
\]
and from the fact that the truncation correction is negligible on the scale
considered, we have
\[
\sqrt{k}\{Q_{n}(x)-A_{n}(x)\}=x^{\beta/\gamma}D_{n,a}^{(\beta)}%
(x)+o_{\mathbf{P}}(1)
\]
in the corresponding weighted integral sense. Furthermore, by regular
variation,
\[
A_{n}(x)^{\beta/p-1}=x^{-(\beta/p-1)/\gamma_{1}}\{1+o_{\mathbf{P}}(1)\},
\]
uniformly in the same weighted integral sense. Since $p=\gamma/\gamma_{1}$,
\[
-\frac{\beta/p-1}{\gamma_{1}}+\frac{\beta}{\gamma}=\frac{1}{\gamma_{1}}.
\]
Therefore,
\[
x^{-1}A_{n}(x)^{\beta/p-1}\sqrt{k}\{Q_{n}(x)-A_{n}(x)\}=x^{-1+1/\gamma_{1}%
}D_{n,a}^{(\beta)}(x)+o_{\mathbf{P}}(1)
\]
in the weighted integral sense. Consequently,
\[
\begin{aligned}
		&\sqrt{k}
		\left[
		\alpha_n\int_{1}^{T_{n,a}}x^{-1}Q_n(x)^{\alpha_n}\,dx
		-
		\frac{\beta}{p}\int_{1}^{T_{n,a}}x^{-1}A_n(x)^{\beta/p}\,dx
		\right] \\
		&\qquad =
		\left(\frac{\beta}{p}\right)^{2}
		\int_{1}^{T_{n,a}}x^{-1+1/\gamma_1}D_{n,a}^{(\beta)}(x)\,dx
		+
		o_{\mathbf P}(1).
	\end{aligned}
\]
Combining this with the definition of $B_{n}$ and
Lemma~\ref{lem:boundary-truncation} gives the desired result.
\end{proof}

\end{document}
